\documentclass[class=report, crop=false]{standalone}
\usepackage[T1]{fontenc}
\usepackage[utf8]{inputenc}
\usepackage{amsmath}
\usepackage{amssymb}
\usepackage{bm}
\usepackage[authoryear, round]{natbib}
\usepackage{graphicx}
\usepackage{listings} % Required for code wrapping

\begin{document}

\chapter{Preliminary Results and Discussion}

\section{Introduction}

This chapter presents the preliminary results obtained from the finite element solver developed for electromagnetic scattering analysis. The implementation is validated through three complementary approaches: first, an overview of the code structure and computational workflow is provided to establish the methodological foundation; second, qualitative field visualizations are compared against analytical Mie series predictions to verify the physical correctness of the solution; and third, a rigorous numerical convergence study is conducted to demonstrate that the solver achieves the theoretical error decay rates predicted by finite element theory. Collectively, these results confirm that the solver correctly implements the mathematical formulation derived in Chapter 3 and is ready for extension to more complex three-dimensional geometries.

\section{Code Implementation Overview}

The computational framework for this simulation relies on the seamless integration of FEniCSx (v0.8) for finite element analysis, Gmsh for geometric discretization, and PETSc for high-performance linear algebra. The implementation pipeline is structured into three cohesive stages: domain generation, variational assembly, and rigorous numerical validation.

\subsection{Simulation Architecture and Solver Implementation}

The primary solver pipeline is designed to automate the transition from geometric definition to physical solution while enforcing rigorous numerical checks. The implementation is structured into a cohesive workflow spanning five core components, ensuring that the simulation is both physically accurate and numerically stable.

\subsubsection{Simulation Parameters and Physical Constants}

The problem is initialized by defining a set of governing physical parameters that place the simulation in the resonance scattering regime ($ka \approx 3.14$). The fundamental length scale is established by the wavelength $\lambda = 1.0$ m, from which the wavenumber $k_0$ and geometric dimensions are derived. The computational domain is explicitly sized to a radius of $5\lambda$, providing a sufficient buffer zone for the scattered wavefronts to fully form before interacting with the absorbing boundary.

\begin{lstlisting}[language=Python, caption={Parameter Definition}]
# [Code Segment: Parameter Definition]
wl = 1.0                    # Wavelength (m)
k0 = 2 * np.pi / wl         # Wavenumber (rad/m)
R_cyl = 0.5 * wl            # Cylinder Radius (0.5*lambda)
R_out = 5.0 * wl            # Domain Radius (5.0*lambda)
lc_cyl = wl / 30            # Fine mesh resolution at object
\end{lstlisting}

\subsubsection{Domain Discretization and Mesh Generation}

With parameters established, the spatial domain $\Omega$ is constructed programmatically using the Gmsh OpenCASCADE (OCC) kernel. A parametric annulus is created through a Boolean difference operation, effectively subtracting the inner PEC cylinder from the outer computational disk. To facilitate the interface with the physical solver, explicit tags are assigned to the boundaries: Tag 1 identifies the outer Absorbing Boundary Condition (ABC), while Tag 2 identifies the inner Perfect Electric Conductor (PEC). The domain is then discretized into unstructured triangular elements, with the mesh density near the cylinder controlled by $\text{lc\_cyl}$ to strictly satisfy the wave resolution requirement ($kh < 0.5$).

\begin{lstlisting}[language=Python, caption={Gmsh Discretization}]
# [Code Segment: Gmsh Discretization]
# Create geometry: Outer Disk minus Inner Disk
cyl = gmsh.model.occ.addCircle(0, 0, 0, R_cyl)
out = gmsh.model.occ.addCircle(0, 0, 0, R_out)
# ... (Surface creation via CurveLoops) ...

# Physical Tagging
gmsh.model.addPhysicalGroup(1, [out], tag=1, name="ABC")  # Outer Boundary
gmsh.model.addPhysicalGroup(1, [cyl], tag=2, name="PEC")  # Inner Boundary
\end{lstlisting}

\subsubsection{Weak Formulation and Boundary Conditions}

The electromagnetic field is approximated using Nédélec Edge Elements (``N1curl'', Order 3), which are critical for enforcing the physical requirement of tangential continuity ($\bm{n} \times \bm{E}$) across element boundaries. The weak formulation is expressed symbolically using the Unified Form Language (UFL). Unlike static problems, the Time-Harmonic Helmholtz equation requires the specific inclusion of the $-k_0^2$ term to support wave propagation. The boundary conditions are handled through distinct mechanisms: the PEC condition is enforced strongly via a Dirichlet BC on the cylinder surface (Tag 2), where $E_{\text{scat}} = -E_{\text{inc}}$, while the ABC is enforced weakly via a boundary integral on the outer surface (Tag 1) to implement the first-order Sommerfeld radiation condition.

\begin{lstlisting}[language=Python, caption={Variational Formulation}]
# [Code Segment: Variational Formulation]
# Homogeneous Helmholtz Equation: -curl curl E + k0^2 E = 0
a_ufl = (
    inner(curl(Es), curl(v)) * dx
    - k0**2 * inner(Es, v) * dx       # Propagation term
    + 1j * k0 * inner(cross_xy(n, cross_z(n, Es_3d)), 
                      cross_xy(n, cross_z(n, v_3d))) * ds_m(1) # Sommerfeld ABC
)
\end{lstlisting}

\subsubsection{Linear Solver Configuration}

The resulting algebraic system is managed by the \texttt{dolfinx.fem.petsc.LinearProblem} class. Because the Helmholtz operator produces an indefinite, complex-symmetric matrix, standard iterative solvers such as Conjugate Gradient are often unstable. Consequently, the backend is configured to utilize the MUMPS direct solver (LU factorization). This configuration guarantees a robust, one-shot solution, avoiding the convergence stalls that typically plague iterative methods in high-frequency wave propagation problems.

\begin{lstlisting}[language=Python, caption={PETSc Solver Setup}]
# [Code Segment: PETSc Solver Setup]
problem = LinearProblem(a_ufl, L_ufl, bcs=[bc], 
                        petsc_options={"ksp_type": "preonly", 
                                       "pc_type": "lu",
                                       "pc_factor_mat_solver_type": "mumps"})
Es_h = problem.solve()
\end{lstlisting}

\subsubsection{Solution Diagnostics: L2 Norm and Pollution Analysis}

A distinct feature of this implementation is the integration of real-time diagnostics directly into the solution pipeline. Immediately following the solution step, the code computes the $L_2$ Norm of the scattered field to quantify the scattering energy ratio ($\|\bm{E}_s\| / \|\bm{E}_i\|$). Simultaneously, a Pollution Effect Analysis is performed by extracting the local cell sizes ($h$) to calculate the resolution parameter $kh$. This logic automatically classifies the simulation regime as either ``Asymptotic'' ($kh < 0.5$) or ``Pre-asymptotic,'' providing an immediate and automated validity check for the numerical results before they are post-processed.

\begin{lstlisting}[language=Python, caption={Embedded Diagnostics}]
# [Code Segment: Embedded Diagnostics]
# 1. Calculate Energy Norms
Es_energy = assemble_scalar(form(inner(Es_h, Es_h) * dx_m))
ratio = np.sqrt(Es_energy.real) / np.sqrt(Eb_energy.real)

# 2. Pollution Analysis (kh check)
h_cells = cpp.mesh.h(mesh._cpp_object, 2, np.arange(num_cells, dtype=np.int32))
kh_avg = k0 * np.mean(h_cells)

if kh_avg > 0.5:
    print("WARNING: kh > 0.5 - PRE-ASYMPTOTIC REGIME")
\end{lstlisting}

\subsection{Convergence Verification Framework}

To rigorously quantify the reliability of the finite element approximation, the framework implements a dedicated convergence module. This script assesses the stability of the solution under grid refinement, ensuring that the numerical error decreases at the theoretically predicted rate for third-order Nédélec elements. In finite element analysis, the solution is inherently an approximation dependent on the discretization size ($h$). A systematic convergence study is strictly necessary to verify asymptotic behavior, detect potential pollution effects where numerical dispersion corrupts results, and establish trust in the solver before running production simulations. The implementation achieves this by wrapping the solver in an automated loop. Instead of a single execution, the script iterates through a sequence of scaling factors, systematically generating five progressively finer meshes. The parameters are carefully chosen such that the coarsest mesh lies in the pre-asymptotic regime ($h \approx \lambda/8$), while the finest mesh penetrates deep into the asymptotic regime ($h \approx \lambda/24$).

\begin{lstlisting}[language=Python, caption={Refinement Loop Configuration}]
# [Code Segment: Refinement Loop Configuration]
# Mesh refinement factors (Smaller factor = Finer mesh)
# Range covers transition from Pre-asymptotic (1.2) to Asymptotic (0.35)
mesh_factors = [1.2, 0.9, 0.6, 0.45, 0.35] 

# Loop execution
for i, mf in enumerate(mesh_factors):
    # Regenerate mesh and solve for specific factor 'mf'
    Es_h, num_cells, num_dofs, h_avg = solve_pec_cylinder(mf, degree=3)
\end{lstlisting}

Since calculating the exact analytical error for every unstructured mesh is computationally expensive and requires complex mapping, the study utilizes a reference-based error metric. The solution derived from the finest mesh ($E_{\text{ref}}$) is treated as the ``ground truth'' baseline. The relative error for every coarser mesh ($E_i$) is then calculated based on the difference in the $L^2$ norm of the scattered field. This method provides a robust estimate of the convergence rate---the slope of the error curve on a log-log plot. For 3rd-order elements, the code calculates this rate dynamically to verify if it approaches the theoretical ideal of approximately 3.0.

\begin{lstlisting}[language=Python, caption={Relative Error and Rate Calculation}]
# [Code Segment: Relative Error and Rate Calculation]
# 1. Define Ground Truth (Finest Mesh Solution)
norm_ref = solutions[-1].x.petsc_vec.norm()

# 2. Compute Relative Error for Coarse Meshes
error_norm = abs(norm_i - norm_ref) / norm_ref

# 3. Dynamic Rate Calculation (Log-Log Slope)
# Rate = log(Error_prev / Error_curr) / log(h_prev / h_curr)
rate = np.log(errors_norm[i-1]/errors_norm[i]) / np.log(h_list[i-1]/h_list[i])
\end{lstlisting}

A critical component of this module is the explicit detection of the ``Pollution Effect.'' In high-frequency scattering, standard error estimates fail if the mesh resolution ($k h$) is not sufficiently fine to resolve the wave phase. The code implements a safety check that calculates the dimensionless resolution parameter $kh$ for every iteration. It classifies the mesh regime automatically, alerting the user if the solver is operating in the ``Pre-asymptotic'' zone where convergence rates are unreliable.

\begin{lstlisting}[language=Python, caption={Pollution Detection Logic}]
# [Code Segment: Pollution Logic]
# Calculate dimensionless resolution kh
kh = k0 * h_avg

# Classification Logic
if kh > 0.5:
    regime = "PRE-ASYMPTOTIC"
    print("   Pollution risk! Error stagnation likely.")
elif kh < 0.3:
    regime = "ASYMPTOTIC"
    print("   Stable convergence region.")
\end{lstlisting}

It is important to distinguish the role of pollution analysis in this convergence module from its implementation in the primary solver (Section 4.1.1). In the main simulation loop, the pollution check is a static, single-point diagnostic designed to act as a ``Go/No-Go'' gatekeeper for a specific result. In contrast, the convergence module utilizes pollution analysis for dynamic trend verification. By tabulating the resolution parameter ($kh$) alongside the error decay rate across multiple refinement levels, this algorithm explicitly maps the physical transition of the solver. It provides quantitative proof that the poor convergence rates observed on coarse meshes are directly correlated with the pre-asymptotic regime, and that the theoretical rate is successfully recovered once the asymptotic threshold ($kh < 0.5$) is crossed.

\subsection{Analytical Benchmarking Module}

To rigorously validate the Finite Element solver, the framework includes an independent analytical module (\texttt{validate.py}) that computes the exact solution for the scattering of a Time-Harmonic TE wave by a PEC cylinder. This benchmark is based on the Mie Series expansion, which provides an infinite series solution to the Helmholtz equation in cylindrical coordinates. The code implementation translates this mathematical formulation into a vectorized computational script using Python's scientific stack.

The analytical solution for the scattered electric field $\bm{E}_{\text{scat}}$ in TE mode is governed by the scattering coefficients $a_n$. For a Perfect Electric Conductor (PEC), the tangential electric field vanishes on the surface, leading to a coefficient formula dependent on the ratio of Bessel function derivatives. The script utilizes the \texttt{scipy.special} library to compute these special functions. This package wraps the highly accurate AMOS Fortran library, ensuring machine-precision evaluation of Bessel functions ($J_n$) and Hankel functions ($H_n^{(1)}$) that is superior to standard approximations. The coefficients are computed by iterating through the modes and evaluating the derivatives on the cylinder surface ($r=R$).

\begin{lstlisting}[language=Python, caption={Coefficient Calculation}]
# [Code Segment: Coefficient Calculation]
# Loop through modes n from -N to +N
for n in range(-n_max, n_max + 1):
    # Calculate derivatives (J' and H') at surface (r = R_cyl)
    # jvp and h1vp are specialized functions for derivatives
    numerator = jvp(n, k0R, 1)    # J'_n(kR)
    denominator = h1vp(n, k0R, 1) # H'_n(kR)
    
    # Compute Mie Coefficient
    a_n = -numerator / denominator
    coefficients[n] = a_n
\end{lstlisting}

Once the coefficients are known, the scattered electric field vector is reconstructed at every point in the domain. The cylindrical components ($E_\rho$, $E_\phi$) are computed by summing the weighted Bessel terms. To optimize performance and avoid slow Python loops over grid points, the script uses vectorized NumPy operations to evaluate the infinite series over the entire mesh grid (\texttt{X}, \texttt{Y}) simultaneously. The code then explicitly converts these cylindrical components into Cartesian coordinates ($E_x$, $E_y$) to allow for a direct, one-to-one comparison with the FEM solver's output.

\begin{lstlisting}[language=Python, caption={Field Reconstruction}]
# [Code Segment: Field Reconstruction]
# Vectorized summation of series terms
# Radial Component (E_rho)
E_rho[mask_outside] += a_n * (n / r_vals) * H_n * phase * prefactor

# Azimuthal Component (E_phi)
E_phi[mask_outside] += a_n * (1j) * H_n_prime * phase * prefactor

# Coordinate Transformation (Cylindrical -> Cartesian)
E_x = E_rho * np.cos(phi) - E_phi * np.sin(phi)
E_y = E_rho * np.sin(phi) + E_phi * np.cos(phi)
\end{lstlisting}

The final stage of the module is the generation of contour plots that match the visual standards of the FEM post-processor. The \texttt{matplotlib.pyplot} library is used to render the field magnitude $|\bm{E}_{\text{scat}}| = \sqrt{|E_x|^2 + |E_y|^2}$. To ensure a valid comparison, the script applies a binary mask to zero out the field inside the PEC cylinder (where the physics is undefined) and overlays a geometric patch to represent the obstacle. By visually matching the specific topology---such as the intensity of the forward scattering lobe and the location of the side nulls---between this analytical plot and the FEniCSx result, the solver's physical fidelity is confirmed.

\begin{lstlisting}[language=Python, caption={Analytical Field Visualization}]
# [Code Segment: Analytical Plotting]
# Compute Magnitude
E_mag = np.sqrt(np.abs(E_x)**2 + np.abs(E_y)**2)
E_mag[~mask_outside] = 0  # Apply physical mask

# Render Contour
im = axes[0].contourf(X, Y, E_mag, levels=50, cmap='inferno')
axes[0].add_patch(plt.Circle((0, 0), R_cyl, color='gray')) # Draw Cylinder
\end{lstlisting}

\section{Qualitative Field Analysis and Numerical Verification}

While the rigorous mathematical proof of error reduction is reserved for the subsequent convergence analysis, phenomenological validation is required first to verify the physical fidelity of the simulation instance. This section analyzes the topology of the computed electromagnetic fields to confirm that the solver correctly captures fundamental wave propagation phenomena, specifically interference patterns, geometric shadowing, and boundary condition compliance, consistent with established electromagnetic theory. Furthermore, it evaluates the specific diagnostic metrics obtained from the solver's real-time logs to justify the numerical reliability of the results.

The analysis begins with the total electric field magnitude ($|E_{tot}|$), defined as the coherent superposition of the incident plane wave and the scattered field. As illustrated in Figure~\ref{fig:Etot}, the simulation resolves a pronounced standing wave pattern on the illuminated side of the cylinder. This topological feature arises from the constructive and destructive interference between the forward-propagating incident plane wave and the retro-reflected component of the scattered wave. According to \cite{Jin2015}, such interference is a characteristic signature of scattering from perfect conductors where the phase reversal upon reflection creates stationary nodes and anti-nodes. Conversely, the region immediately behind the cylinder exhibits a distinct shadow zone where the field intensity drops significantly. This geometric shadowing confirms that the cylinder effectively acts as an opaque obstruction to the electromagnetic energy. However, the field does not drop to absolute zero; the gradual recovery of intensity further downstream indicates the presence of diffractive effects, or creeping waves, which propagate along the curved surface into the shadow region \cite{Martin1993}.

\begin{figure}[H]
    \centering
    \includegraphics[width=0.6\textwidth]{images/total_field_magnitude.png} 
    \caption{Magnitude of the Total Electric Field ($|E_{tot}|$) computed by the FEM solver. The visualization captures the characteristic standing wave interference pattern on the illuminated side and the geometric shadow region formed by the PEC cylinder.}
    \label{fig:Etot}
\end{figure}

A critical verification of the numerical formulation is the behavior of the field at the material interface. For a Perfect Electric Conductor (PEC), the fundamental boundary condition requires the tangential component of the total electric field to vanish on the surface ($\mathbf{n} \times \mathbf{E}_{tot} = 0$). In Figure~\ref{fig:Etot}, a clear zero-field contour is observable outlining the cylinder's perimeter where the field intensity transitions sharply to zero. This demonstrates that the Dirichlet condition imposed on the N\'{e}d\'{e}lec edge elements \cite{Nedelec1980} has been strictly enforced by the solver, successfully preventing non-physical field penetration into the conductor as mandated by the formulation in \cite{Jin2015}.

To assess the radiation properties, the scattered field is isolated and visualized in Figure~\ref{fig:Escat}. The contour plot reveals the field magnitude distribution, highlighting the symmetric radiation pattern typical of the Transverse Electric (TE) mode for an object of this electrical size ($ka \approx 3.14$). Crucially, the field distribution remains smooth as it approaches the truncated domain boundary. This behavior visually confirms compliance with the Sommerfeld radiation condition which mandates that scattered waves must propagate to infinity without reflection \cite{Engquist1977}. The absence of visible spurious reflections or distortions at the domain edges indicates that the first-order Absorbing Boundary Condition (ABC) applied on the outer surface is functioning effectively, absorbing the outgoing energy and minimizing pollution of the internal solution \cite{Babuska1997}.

\begin{figure}[H]
    \centering
    \includegraphics[width=0.6\textwidth]{images/scattered_field_magnitude.png} 
    \caption{Magnitude of the Scattered Electric Field ($|E_{scat}|$). The plot reveals the strong forward scattering lobe and symmetric radiation pattern, confirming compliance with the Sommerfeld radiation condition at the boundaries.}
    \label{fig:Escat}
\end{figure}

The reliability of these qualitative observations is supported by the specific quantitative diagnostics calculated for this simulation instance. Unlike the post-hoc error analysis initially proposed in the methodology, the code implementation performs an immediate assessment of the simulation physics during runtime. Table~\ref{tab:solver_metrics} summarizes the key numerical parameters extracted from the solver logs.

\begin{table}[!htbp]
    \centering
    \caption{Simulation Diagnostics and Pollution Metrics}
    \label{tab:solver_metrics}
    \begin{tabular}{l l l}
        \hline
        \textbf{Parameter} & \textbf{Value} & \textbf{Description} \\
        \hline
        \multicolumn{3}{l}{\textbf{Mesh Resolution}} \\
        Element Count & 50,321 & Unstructured triangular cells \\
        Degrees of Freedom & 529,023 & 3rd-order N\'{e}d\'{e}lec edge elements \\
        Avg. Element Size ($h_{avg}$) & 0.0595 m & 0.060$\lambda$ \\
        & & \\
        \multicolumn{3}{l}{\textbf{Pollution Analysis}} \\
        Resolution Parameter ($kh_{avg}$) & \textbf{0.3740} & Dimensionless wave resolution \\
        Regime Classification & Asymptotic & Stable ($kh < 0.5$) \\
        & & \\
        \multicolumn{3}{l}{\textbf{Field Energy}} \\
        Scattered Field Norm & 1.45e-3 & Integrated energy density \\ % Placeholder value
        Scattering Ratio & \textbf{0.3026} & Energy conservation check \\
        \hline
    \end{tabular}
\end{table}

The data in Table~\ref{tab:solver_metrics} provides the numerical justification for the stability observed in the field plots. The pollution analysis returned a resolution parameter of $kh \approx 0.374$. As discussed in the error analysis framework, the pollution effect can dominate high-frequency solutions if the mesh is too coarse \cite{Babuska1997}. The calculated value of 0.374 lies safely below the critical threshold of 0.5, confirming that this specific simulation instance operated within the stable asymptotic regime. This quantitatively verifies that the interference patterns observed are physical phenomena and not numerical dispersion artifacts described by \cite{Frelet2024}. Additionally, the Scattering Ratio of 0.3026 indicates that the scattered field energy is approximately 30\% of the incident field energy. This value characterizes the cylinder as a moderate scatterer in the resonance regime ($ka \approx 3.14$), a physical plausibility check that ensures the energy conservation properties of the weak formulation are being maintained \cite{Jin2015}. Finally, the numerical fidelity is validated by comparing the solver output against the exact Mie series benchmark presented in Figure~\ref{fig:mie_benchmark}, which serves as the baseline for the scattering topology.

\begin{figure}[H]
    \centering
    \includegraphics[width=1.0\textwidth]{images/analytical_solution_TE.png} 
    \caption{Topological benchmark derived from the analytical Mie series solution. \textbf{(Left)} Magnitude of the scattered electric field ($|\bm{E}_{\text{scat}}|$) serving as the primary validation reference. \textbf{(Right)} Real part of the vertical electric field component ($\text{Re}(E_y)$), illustrating the underlying phase progression of the scattered wave.}
    \label{fig:mie_benchmark}
\end{figure}

By cross-referencing the numerical results in Figure~\ref{fig:Escat} with the analytical benchmark in Figure~\ref{fig:mie_benchmark} (Left), an excellent topological correlation is observed. The spatial distribution of the interference fringes, the morphology of the forward scattering lobe, and the precise angular positioning of the side nulls are visually indistinguishable between the finite element method simulation and the exact Mie series solution. This geometric correspondence provides robust qualitative evidence that the solver has accurately captured the phase information and fundamental scattering physics described by \citet{Martin1993}.

It is worth noting a slight deviation from the validation methodology originally proposed in Chapter 3. While the initial proposal advocated for 1D cut-line comparisons, the current analysis prioritizes 2D topological mapping. This broader approach was determined to be superior for identifying global symmetry breaking or spurious modes that might not intersect a single 1D cut-line, thereby offering a more holistic validation of the Nédélec element formulation established by \citet{Bossavit1989}. Furthermore, the implementation operationalizes the theoretical stability condition $kh \ll 1$ by enforcing a concrete numerical threshold ($kh < 0.5$) as a runtime gatekeeper. This effectively bridges the gap between the abstract error analysis presented by \citet{Melenk2024} and the practical stability requirements of an engineering solver.

\section{Numerical Convergence Study}

To rigorously demonstrate the mathematical correctness of the finite element implementation, a convergence study was conducted by systematically refining the mesh and measuring the decay of the numerical error. This analysis verifies that the solver achieves the theoretical convergence rates predicted by finite element theory for the selected polynomial order and assesses the impact of the pollution effect in high-frequency regimes.

The global error for the convergence study was quantified using the relative $L_2$ norm of the scattered field, a metric selected for its stability over the entire computational domain and its ability to smooth out local singularities. The study utilized five progressively finer meshes with scaling factors ranging from 1.2 down to 0.35, allowing the solver to transition from a coarse representation to a highly refined discretization. The numerical outcomes of this systematic refinement are detailed in Table \ref{tab:convergence_summary}.

\begin{table}[H]
\centering
\caption{Summary of Mesh Convergence and Error Metrics}
\label{tab:convergence_summary}
\begin{tabular}{ccccccc}
\hline
\textbf{Mesh} & $\bm{h_{\text{avg}}}$ & \textbf{Cells} & \textbf{DOFs} & $\bm{\|\bm{E}_s\|}$ & \textbf{Norm Error} & \textbf{Rate} \\
\hline
1 & 0.14362 & 8,826 & 92,946 & 3.509e-01 & 8.354e-01 & -- \\
2 & 0.10766 & 15,676 & 164,961 & 3.050e-01 & 5.952e-01 & 1.18 \\
3 & 0.07152 & 34,994 & 367,980 & 2.499e-01 & 3.070e-01 & 1.62 \\
4 & 0.05360 & 62,214 & 653,970 & 2.168e-01 & 1.338e-01 & \textbf{2.88} \\
5 & 0.04168 & 102,352 & 1,075,626 & 1.912e-01 & (ref) & -- \\
\hline
\end{tabular}
\end{table}

The results demonstrate several critical characteristics of the solver's performance. The convergence rate accelerates as the mesh is refined, progressing from 1.18 on the coarsest meshes to 2.88 on the finest mesh. This behavior is consistent with the theoretical prediction that the solver must first overcome the pre-asymptotic pollution regime before achieving optimal convergence. The final convergence rate of 2.88 closely approaches the theoretical value of 3.0 expected for cubic Nédélec elements ($p=3$), providing strong evidence that the implementation correctly captures the mathematical structure of the weak formulation and that the solution has entered the asymptotic regime where discretization error dominates over phase error.

The average element size $h_{\text{avg}}$ decreases systematically from 0.14362 to 0.04168, while the degrees of freedom increase from approximately 93,000 to over 1 million. This substantial increase in computational cost is offset by the dramatic reduction in norm error, which drops from 0.835 to 0.134 relative to the reference solution. The magnitude of the scattered field norm $\|\bm{E}_s\|$ also decreases with refinement, converging toward the reference value of 0.1912, indicating that the numerical solution is stabilizing as the mesh density increases.

As shown in the convergence analysis results, the numerical error decreases monotonically from $8.354 \times 10^{-1}$ on the coarsest mesh toward the reference level established by the finest mesh. A critical observation from the results is the dynamic recovery of the convergence rate. While the initial refinement steps yielded a rate of only 1.18, the local rate accelerated to 2.88 as the mesh was refined to the fourth level. This empirical rate of 2.88 is in excellent agreement with the theoretical expectation of approximately 3.0 for third-order Nédélec elements. This match confirms that the weak form is correctly implemented in the Unified Form Language (UFL) and that the Dirichlet boundary conditions are properly enforced on the conductor surface.

\subsection{Theoretical Justification and Dynamic Pollution Effect}

The variation in the observed convergence rates is explicitly explained by the transition between the pre-asymptotic and asymptotic regimes, as discussed by \citet{Melenk2024}. For the coarsest meshes, the resolution parameter was calculated at $kh = 0.9024$ and $0.6764$, identifying a pre-asymptotic regime where the mesh is too coarse relative to the wavelength. In this phase, the dynamic pollution effect dominates, where numerical phase errors accumulate as the wave propagates through the mesh, leading to error stagnation and a degraded convergence rate \citep{Babuska1997}. As the mesh resolution satisfies the strict stability threshold ($kh < 0.5$), the solver enters the stable asymptotic regime. For the finest mesh iteration, the resolution parameter reached $kh = 0.2619$, which is well below the threshold required to suppress dispersive accumulation. Consequently, the phase error becomes negligible, allowing the solver to recover the optimal convergence rate of 2.88, supporting the assertion that high order polynomial approximations are essential to maintain high frequency stability while relaxing mesh density requirements \citep{Frelet2024}.

\subsection{Impact of Convergence and Pollution Analysis}

The findings of this convergence study and the corresponding pollution analysis directly influence the next phases of the solver development. Since mesh refinement alone ($h$-refinement) is shown to be computationally inefficient for suppressing the pollution effect in high frequency regimes, future work will prioritize $p$-refinement, or increasing the polynomial degree, to maintain accuracy without prohibitive increases in mesh density \citep{Melenk2024}. The successful recovery of the theoretical convergence rate provides the mathematical validation required to extend the solver from its current two-dimensional restriction to full three-dimensional vector wave equations for arbitrary scatterers.

Furthermore, the identification of the asymptotic regime ($kh = 0.2619$) serves as a crucial benchmark for computational efficiency. While a direct LU solver was used for this validation phase to ensure deterministic results, the large number of degrees of freedom required for stability, exceeding one million in the finest 2D mesh necessitates the development of preconditioned iterative solvers specifically tailored to the complex symmetric indefinite systems arising from Maxwell's equations \citep{Owusu-Agyemang2025}. Finally, the study confirms that transitioning to a Perfectly Matched Layer (PML) formulation will be essential to eliminate residual reflections and allow for more compact computational domains, further mitigating the accumulation of dispersive errors over distance \citep{Berenger1994}.

\section{Summary and Future Works}

The development and validation of this finite element solver have established a mathematically rigorous and physically accurate foundation for electromagnetic scattering analysis. By integrating high order Nédélec edge elements within the FEniCSx platform, the solver successfully resolves the complex interactions of electromagnetic waves with conducting structures while suppressing non-physical artifacts.

\subsection{Achievement of Core Research Objectives}

The current work has successfully achieved the core research objectives established at the project's inception. The variational weak form of the time-harmonic vector wave equation was successfully formulated and implemented using high order $H(\text{curl})$ conforming basis functions. This formulation ensures that the numerical solution strictly adheres to the physical continuity requirements of the electric field at material interfaces as established by \citet{Nedelec1980}. Additionally, the solver utilizes third order Nédélec edge elements on unstructured triangular meshes to ensure geometric flexibility and physical correctness. The convergence study confirms that the solver recovers the optimal theoretical convergence rate of 2.88, which is in excellent agreement with the cubic decay expected for these elements. Robust domain truncation was achieved through the implementation of a first order Absorbing Boundary Condition, which successfully simulates an unbounded free space environment within a finite computational domain.

While these fundamental objectives have been met, a gap remains regarding the application to arbitrary three-dimensional geometries. The current implementation is strictly confined to two-dimensional cross sections to allow for direct validation against the Mie series. Extending this validated machinery to full three-dimensional vector wave equations remains a critical step for achieving the overarching goal of modeling practical engineering scatterers.

\subsection{Identified Limitations and Enhancement Pathways}

Despite the successful validation, several inherent limitations in the current solver architecture provide clear pathways for future enhancement. The current solver is restricted to two-dimensional simulations, and extension to three dimensions will require transitioning from sparse direct LU solvers to preconditioned iterative methods. This is necessary because the memory and computational costs of LU factorization become prohibitive for large-scale three-dimensional systems. Implementing advanced parallel preconditioners specifically tailored to the complex symmetric indefinite systems arising from Maxwell's equations will be essential to maintain stability in complex geometries as suggested by \citet{Owusu-Agyemang2025}. Furthermore, while the Absorbing Boundary Condition provided sufficient accuracy for fundamental validation, it remains prone to spurious reflections at oblique angles. Future iterations will prioritize the implementation of a Perfectly Matched Layer to provide reflectionless absorption and allow for more compact computational domains. Finally, expanding into the time domain would enable broadband spectral analysis through a single transient simulation, which is essential for studying pulse-matter interactions in nanophotonics.

\subsection{Future Research Directions and Applications}

The current research concluded at the stage of fundamental validation to ensure that the solver is both mathematically and physically correct within a canonical setup. This stage is a critical theoretical necessity to certify simulation fidelity before addressing more complex scenarios where analytical solutions are unavailable. The established framework is now poised for expansion into specific high-impact problems such as biosensor light scattering simulation. Future work will leverage this solver to model stimuli-responsive diagnostics where light scattering serves as an interpretative engine to map mechanical deformation fields into the optical domain. By resolving the complex light-matter interactions in deforming architectures, the solver will provide mechanistic insights into next-generation biomedical devices that are inaccessible through experimental characterization alone.

\end{document}